\section{Discussion}
\label{sec:discussion}

Three conclusions are supported consistently by the analytical results of Section~\ref{sec:analysis} and the numerical layers of Sections~\ref{sec:response}--\ref{sec:benchmark}.

First, \emph{structural identifiability and practical discrimination are different questions}. The Caputo transfer carries an exact $O(\omega^{-\alpha})$ signature (Theorem~\ref{thm:T4}); finite latent models approximate the same response to any tolerance on $[0,T]$ (Corollary~\ref{cor:closure}); and the testing bounds quantify the gap --- $P_e^*\ge0.498$ at $m=32$ under the pulse protocol (Table~\ref{tab:state-approx}). A flexible latent rival therefore remains problematic long after the model classes are proved analytically distinct.

Second, \emph{delay and fractional memory should not be treated as mutually exclusive}. The simulations of Section~\ref{sec:fd} trace a continuous bridge between the Caputo limit ($\tau=0$) and the DDE limit ($\alpha=1$): delay governs predator-response timing, fractional order the relaxation structure, and interior cells sit far from both limits simultaneously ($d_{\rm Caputo}=0.28$, $d_{\rm DDE}=0.69$ at $(\alpha,\tau)=(0.95,0.50)$; Table~\ref{tab:fd-bridge-cells}). Fitting the two endpoint models separately cannot reveal this; the composite stress test can.

Third, \emph{within classical waveform families the most informative perturbation is not the safest}, by a wide margin: the accuracy-leading inputs cross the Allee threshold at rates $0.357$--$1.000$ in the four-class benchmark (Table~\ref{tab:benchmark-safety}) and in $3/9$ to $9/9$ cells of the fractional-delay grid (Table~\ref{tab:fd-design-ranking}), two independent measurements of the same effect. Crossing $A$ is not a numerical artifact: it changes the basin structure, and the procedure ceases to be a local identification experiment at that point.

Fourth, and this qualifies the third conclusion decisively, \emph{that trade-off belongs to the parameterisation and not to the ecology}. Neither reducing amplitude (Section~\ref{subsec:amplitude}) nor accepting the classical frontier is necessary: searching piecewise-constant inputs at the same amplitude budget under a hard margin constraint produces designs that are simultaneously safe and more informative than every classical design, safe or not --- macro-accuracy $0.803$ with zero crossings against $0.514$ for the only safe classical family and $0.763$ for the best classical design overall (Table~\ref{tab:safe-search}). The frontier of Table~\ref{tab:benchmark-safety} is therefore a statement about six waveform families, and reporting it as a property of safe ecological identification would be an error. What remains binding is the latent-complexity obstruction of Corollary~\ref{cor:closure}, Theorem~\ref{thm:T20} and Tables~\ref{tab:state-approx}--\ref{tab:state-approx-large}: the induced testing floor reaches $0.49999$ by $m=128$, independently of the decision rule. The benchmark's low recall against the finite-latent rival ($0.475$ for the leading safe design, and at or below chance for the class overall) is consistent with this, but BIC's complexity penalty is a competing explanation that the present design does not separate, so the obstruction rests on Section~\ref{subsec:state} and not on the recall. The operative constraint on identifying ecological memory is the declared rival-complexity budget, not the safety margin.

\paragraph{Limitations.}
Theorem~\ref{thm:T23} is linearized and specific to the prey channel; a full-state, hierarchy-wide safety closure remains open. The waveform search of Section~\ref{subsec:safe-search} establishes existence, not global optimality: no certified global bound over the piecewise-constant class is claimed, the search was run at a single amplitude and a single $\alpha$, and its safety constraint is active at the declared margin. Safety of the leading designs is verified a posteriori by mesh refinement together with a rigorous inter-node bound from the Caputo modulus of continuity; this decides every trajectory tested but is not an interval enclosure of the solver error, and a validated enclosure was obtained for only four of the eight headline trajectories. The fractional-delay study is numerical only, with constant equilibrium history, one delayed trophic pathway, and the targeted grids of Table~\ref{tab:experiment-grid}. The benchmark decides by BIC, not Bayesian evidence; adaptive design is specified (Supplementary Material) but not executed. This ledger separates proved results from computational evidence.

\paragraph{Relation to the previously submitted companion paper.}
The shared material is the strong-Allee backbone with its locked parameterization, from Ref.~\cite{P01}; nothing else is shared. Ref.~\cite{P01} concerns certification: computer-assisted verification of Caputo equilibria, the Matignon boundary $\alpha^*(A)$, extinction/boundedness, and the diagnostic failure of a forward-Euler surrogate. The present work concerns discrimination. Theorems~\ref{thm:T4}, \ref{thm:T20} and~\ref{thm:T23}, Lemma~\ref{lem:soe}, the six-family design study, the 1512-cell benchmark, and the fractional-delay simulations appear nowhere in Ref.~\cite{P01}. Each inherited stability fact enters only to select an admissible operating regime, and is attributed where it enters.

\paragraph{What is genuinely added by the present paper.}
Backbone reuse is an experimental-design decision. With the ecology frozen, differences between simulated observations are attributable to the memory law or to the experiment, never to an altered population model. The paper starts where~\cite{P01} ends: a certified operating point leaves open whether an observed transient identifies a fractional mechanism uniquely, whether a delayed or latent model can reproduce it (Corollary~\ref{cor:closure} says yes, at sufficient $m$), and which perturbations separate the alternatives within the safe regime. These questions require the mathematics of Section~\ref{sec:analysis} and the computations of Sections~\ref{sec:numerics}--\ref{sec:benchmark}, none of which appears in~\cite{P01}.

The figure sets reflect the same division: the centerpiece of~\cite{P01} is a certified $(A,\alpha)$ stability atlas, whereas Figures~\ref{fig:transfer-mag}--\ref{fig:safety-tradeoff} here concern response separation, approximation error, waveform design, BIC confusion, fractional-delay dynamics, and the safety--informativeness frontier. Only the model schematic (Figure~\ref{fig:schematic}) orients the reader in the common ecology; no certificate atlas or certificate count from~\cite{P01} is used as evidence.

\paragraph{Relation to broader prior work.}
Safe information limits and safe active model discrimination exist as general theories~\cite{SAFE26A,SAFE26B}; no novelty is claimed for safe experimental design per se. The contribution is the specific quantitative interaction: finite-horizon fractional approximation (Lemma~\ref{lem:soe}, a borrowed construction), bounded latent complexity, strong-Allee geometry, and the combined fractional-delay stress test in a single frame.

\paragraph{Implications for experimental identification.}
Model selection here is a \emph{controlled intervention problem}, not passive curve fitting. Under weak excitation, BIC collapses to the ODE explanation (recall $0.972$, precision $0.227$); stronger excitation separates the mechanisms and simultaneously raises the crossing rate. The fractional-delay study adds a distinct warning: both endpoint models can be individually identifiable while an interior composite occupies a response region neither represents (Table~\ref{tab:fd-bridge-cells}), so designs should test composite alternatives, not only the endpoints.

The operative objective is not ``maximize information'' but maximize information \emph{subject to} a trajectory-level safety margin and a declared rival-complexity budget --- and Section~\ref{subsec:safe-search} shows that this constrained problem should actually be solved rather than approximated by choosing among textbook waveforms, since the loss from not solving it is $0.29$ in macro-accuracy. The second constraint is forced by Corollary~\ref{cor:closure}: an unrestricted latent hierarchy reproduces fractional memory on any finite horizon, so the allowed latent dimension $m$, delay range, forcing bandwidth, and observation horizon $T$ belong in the scientific specification of the experiment, alongside organism and protocol.
